\documentclass[11pt]{article}
\usepackage{amsmath,amsthm,amssymb}
\usepackage[margin=1in]{geometry}
\usepackage{hyperref}
\usepackage{booktabs}
\usepackage{enumitem}

\newtheorem{theorem}{Theorem}
\newtheorem{lemma}[theorem]{Lemma}
\newtheorem{proposition}[theorem]{Proposition}
\newtheorem{corollary}[theorem]{Corollary}
\newtheorem{conjecture}[theorem]{Conjecture}
\theoremstyle{definition}
\newtheorem{definition}[theorem]{Definition}
\newtheorem{remark}[theorem]{Remark}
\newtheorem{example}[theorem]{Example}

\newcommand{\Hq}{H_q}
\newcommand{\Sn}{S_n}
\newcommand{\Vlam}{V_\lambda}
\newcommand{\Bdec}{B_{\ell+1}^{(\lambda)}}
\newcommand{\Om}{\Omega}
\newcommand{\hatl}{\widehat\lambda}
\newcommand{\decl}{\lambda^{(\ge 2)}}

\title{Extended sign-kill for $\hatl=(2,2)$: \\
  the complete $E_j^-$ chain at all $q\ge 2$}
\author{Clio}
\date{14 May 2026 (very late prove session, second pass)}

\begin{document}
\maketitle

\begin{abstract}
Let $\lambda = (2,2,1^q) \vdash n = q+4$ with $q \ge 2$, $\ell = q+2$,
and $B = \Bdec\Vlam := R'_{\ell+1}\cdots R'_n\,\Vlam$.  We prove the
extended sign-kill conjecture for this family:
\[
   B \;\subseteq\; \bigcap_{j=1}^{q-1} E_j^-\!\mid_{\Vlam}
   \quad\text{and}\quad
   B \not\subseteq E_j^-\!\mid_{\Vlam} \text{ for } j \ge q.
\]
Combined with the universal $B \subseteq E_\ell^+$ (Theorem~A,
May-13 evening), this completely identifies the $T_j$-eigenspace
containments of $B$ for the family.

The proof is short.  The May-13 paper
\texttt{2026-05-13-non-hook-22-1n.tex} established $\dim B = 1$ and
identified the explicit $4$-SYT support of the spanning vector $b$.
We observe that this support immediately forces the column-$1$
entries $1, 2, \ldots, n-4$ to occupy column-$1$ rows $1, 2, \ldots,
n-4$ in every support SYT.  The Hoefsmit same-column rule gives
$T_j v_T = -v_T$ for each support SYT and $1 \le j \le n-5 = q-1$,
hence $T_j b = -b$.  Sharpness for $j = q, q+1, \ldots$ follows
from the non-vanishing of the four support coefficients (also from
May-13) together with the off-diagonal Hoefsmit blocks at axial
distance $\ge 2$.

This closes the $\hatl = (2,2)$ row of the May-14 extended sign-kill
conjecture table, in both directions.
\end{abstract}

\section{Setup}\label{sec:setup}

We use the conventions of the May-13--May-14 series.  $\Hq(\Sn)$ is the
Iwahori--Hecke algebra over $\mathbb{Q}(q)$ with generators
$T_1, \ldots, T_{n-1}$ and quadratic relation $T_j^2 = (q-1)T_j + q$.
$\Vlam$ is the Specht module in the Hoefsmit seminormal basis
$\{v_T : T \in \mathrm{SYT}(\lambda)\}$.

For $1 \le j \le n-1$, set $S_j := T_j + 1$.  The Hecke relation
factors as $S_j(T_j - q) = 0$, equivalently $T_j S_j = qS_j = S_j T_j$.
Write
\[
   E_j^+ \;:=\; \ker(T_j - q)\mid_{\Vlam} \;=\; \mathrm{Im}(S_j\mid_{\Vlam}),
   \qquad
   E_j^- \;:=\; \ker(S_j)\mid_{\Vlam}.
\]
Since the $T_j$ acts with minimal polynomial $(T-q)(T+1)$ and the two
eigenvalues are distinct (for generic $q$), $\Vlam = E_j^+ \oplus E_j^-$.

For $k \ge 1$, let $R'_k := S_{k-1}\,S_{k-2}\cdots S_1$.  For
$\lambda \vdash n$ with $\ell = \ell(\lambda)$, set
\[
   \Om \;=\; \Om^{(\lambda)}_\ell \;:=\; R'_{\ell+1}R'_{\ell+2}\cdots R'_n
   \;\in\;\Hq,
   \qquad
   B \;:=\; \Om \cdot \Vlam.
\]

For $\lambda = (2,2,1^q)$, $n = q+4$ and $\ell = q+2$, so
$\Om = R'_{q+3}R'_{q+4}$ has two $R'$-blocks.  By the May-13
multi-corner-dimension paper, $\dim B = f^{(1,1)} = 1$ for $q \ge 0$
(stable regime starts at $q = 0$ here).

\paragraph{SYT labelling for $\lambda = (2,2,1^q)$.}  An SYT $T$ of
$\lambda$ has cell $(1,1)$ always equal to $1$ and column~$2$ consisting
of two cells $(1,2)$ and $(2,2)$.  Setting $a := T(1,2)$ and
$c := T(2,2)$, we have $2 \le a < c \le n$ and $(a,c) \ne (2,3)$
(the latter forced by row-$2$ increasing).  Conversely, every such
pair $(a,c)$ yields a unique SYT, denoted $T^{(a,c)}$, by filling
column~$1$ at rows $1, 2, \ldots, n-2$ with the elements of
$\{1\} \cup \{2,\ldots,n\}\setminus\{a,c\}$ in increasing order.
Total count: $\binom{n-1}{2} - 1 = f^\lambda$.

\section{Input from prior papers}\label{sec:input}

\begin{theorem}[Support of $B$, May-13]\label{thm:may13-support}
For $\lambda = (2,2,1^q)$ with $q \ge 2$ (i.e.\ $n \ge 6$),
\[
   B = \mathrm{span}(b), \qquad
   b \;=\; \sum_{(a,c) \in \mathcal{S}_n} C_{(a,c)}\,v_{T^{(a,c)}},
\]
where
\[
   \mathcal{S}_n \;:=\; \bigl\{(n-3,n-2),\,(n-3,n-1),\,(n-2,n),\,(n-1,n)\bigr\}
\]
and all four coefficients $C_{(a,c)} \in \mathbb{Q}(q)$ are nonzero.
\end{theorem}

\begin{proof}
This is precisely Theorem~1 (combined with Proposition~5.4 on the
terminal step) of \texttt{2026-05-13-non-hook-22-1n.tex}.
\qed
\end{proof}

\begin{theorem}[Universal $T_\ell$ scalar, May-13 evening]\label{thm:T-ell}
For any partition $\lambda \vdash n$ with $1 \le \ell = \ell(\lambda) \le n-1$,
$T_\ell \cdot \Om^{(\lambda)} = q \cdot \Om^{(\lambda)}$ in $\Hq$.  In
particular, $B \subseteq E_\ell^+\mid_{\Vlam}$.
\end{theorem}

\section{The extended sign-kill theorem}\label{sec:main}

\begin{theorem}[Main]\label{thm:main}
For $\lambda = (2,2,1^q)$ with $q \ge 2$, let $\ell = q+2$, $n = q+4$.
The image $B = \Om\cdot\Vlam$ satisfies
\[
   B \;\subseteq\; E_\ell^+\;\cap\; \bigcap_{j=1}^{q-1} E_j^-\!\mid_{\Vlam},
\]
and this is sharp:
\[
   B \not\subseteq E_j^- \text{ for } j \in \{q, q+1, \ldots, n-1\},
   \qquad
   B \not\subseteq E_j^+ \text{ for } j \in \{1,\ldots,n-1\}\setminus\{\ell\}.
\]
\end{theorem}

We prove the four containments and the four non-containments
separately.

\subsection{$E_\ell^+$ containment}

Immediate from Theorem~\ref{thm:T-ell}.

\subsection{$E_j^-$ containment for $1 \le j \le q-1$}

The key combinatorial observation:

\begin{lemma}[Column~$1$ structure of support SYTs]\label{lem:top-col1}
For $(a, c) \in \mathcal{S}_n$ with $a < c$, column~$1$ of $T^{(a,c)}$
holds the elements of $\{1, 2, \ldots, n\} \setminus \{a, c\}$ at rows
$1, 2, \ldots, n-2$ in increasing order; in particular
$T^{(a,c)}(i, 1) = i$ if and only if $i \le a-1$.
\end{lemma}

\begin{proof}
By definition of $T^{(a,c)}$.  Since $a < c$, the sorted column-$1$
sequence is
\[
   1,\, 2,\, \ldots,\, a-1,\;\; a+1,\, a+2,\, \ldots,\, c-1,\;\; c+1,\, c+2,\, \ldots,\, n.
\]
The $i$-th element (occupying row $i$ of column~$1$) equals $i$ iff
$i \le a-1$.
\qed
\end{proof}

\begin{remark}[The three bounds we will use]\label{rem:bounds}
By Lemma~\ref{lem:top-col1}, the column-$1$ ``identity prefix'' has
length $a - 1$ in $T^{(a,c)}$.  We will repeatedly use:
\begin{itemize}[topsep=0.2em,itemsep=0.1em]
\item All four members of $\mathcal{S}_n$ have $a \ge n - 3$,
  giving $T^{(a,c)}(i,1) = i$ for $i \le n - 4$.
\item $(n-2, n)$ and $(n-1, n)$ have $a \ge n - 2$, giving the
  additional $T^{(a,c)}(n-3, 1) = n-3$.
\item $(n-1, n)$ has $a = n - 1$, giving also $T^{(a,c)}(n-2, 1) = n-2$.
\end{itemize}
\end{remark}

\begin{proposition}[$E_j^-$ containment]\label{prop:Ej-minus}
For $1 \le j \le q - 1$ and every $(a, c) \in \mathcal{S}_n$,
$T_j\,v_{T^{(a,c)}} = -v_{T^{(a,c)}}$.  Consequently
$T_j\,b = -b$, i.e., $b \in E_j^-$.
\end{proposition}

\begin{proof}
Let $T := T^{(a,c)}$.  For $1 \le j \le q-1 = n-5$, both $j$ and $j+1$
lie in $\{1, 2, \ldots, n-4\}$.  By Lemma~\ref{lem:top-col1}, the
entries $j$ and $j+1$ are at cells $(j, 1)$ and $(j+1, 1)$ of $T$;
in particular they occupy the same column (column~$1$) at adjacent
rows.  By the Hoefsmit seminormal action (same-column case),
$T_j v_T = -v_T$.

By linearity over $(a, c) \in \mathcal{S}_n$ in
$b = \sum_{(a,c)} C_{(a,c)} v_{T^{(a,c)}}$ (Theorem~\ref{thm:may13-support}),
$T_j b = -b$.
\qed
\end{proof}

\subsection{$E_j^-$ non-containment for $q \le j \le n-1$}

The argument uses non-vanishing of the four support coefficients
together with the existence of an off-diagonal Hoefsmit block.

\begin{lemma}[Off-diagonal action at $j = q$]\label{lem:j=q-offdiag}
For $j = q$ (so $j + 1 = q + 1$) and the support SYT
$T^* := T^{(n-3, n-2)} = T^{(q+1, q+2)}$, the entries $q$ and $q+1$
sit at cells $(q, 1)$ and $(1, 2)$ of $T^*$ respectively (different
row and different column).  The corresponding axial distance is
$d = c(1,2) - c(q,1) = 1 - (1-q) = q$.  In particular, the
Hoefsmit $T_q$-action on the $2$-block
$\{v_{T^*}, v_{s_q T^*}\}$ is non-degenerate over $\mathbb{Q}(q)$,
i.e., $T_q v_{T^*} = \alpha\,v_{T^*} + \beta\,v_{s_q T^*}$
with $\beta \ne 0$.

Here $s_q T^* = T^{(q, q+2)}$ (swap entries $q$ and $q+1$, which
exchanges $T^*(q,1) = q$ with $T^*(1,2) = q+1$); this is a valid
SYT of $\lambda$ and is \emph{not} an element of $\mathcal{S}_n$.
\end{lemma}

\begin{proof}
Lemma~\ref{lem:top-col1} gives $T^*(q, 1) = q$.  By definition of
$T^* = T^{(q+1, q+2)}$, $T^*(1, 2) = q+1$.  So $j = q$ at $(q,1)$
(content $1 - q$) and $j+1 = q+1$ at $(1,2)$ (content $1$).
Different row and column, axial distance
$d = 1 - (1-q) = q$.  For $q \ge 2$ ($\Leftrightarrow$ $|d| \ge 2$),
the Hoefsmit ratio
$\tau(d) = (q^{|d|}-1)/(q^d \cdot (q-1))\cdot(\cdot)$ (the precise
formula does not matter) gives $\beta \ne 0$ over $\mathbb{Q}(q)$.
Concretely, the off-diagonal entry of the standard Hoefsmit
$2$-block at axial distance $d$ is a nonzero rational function in
$q$ for all $|d| \ge 2$.

Swapping $q$ and $q+1$ in $T^*$ yields the SYT $T^{(q, q+2)}$ (with
column-$2$ entries $q, q+2$).  Since $q < q+1$ and $(q, q+2) \ne
\mathcal{S}_n$ (which only contains pairs with both entries $\ge n-3 = q+1$),
this swapped SYT lies outside the support of $b$.
\qed
\end{proof}

\begin{proposition}[$B \not\subseteq E_q^-$]\label{prop:not-Eq-minus}
$T_q b \ne -b$, hence $b \notin E_q^-$.
\end{proposition}

\begin{proof}
By Theorem~\ref{thm:may13-support}, $b = \sum_{(a,c)\in\mathcal{S}_n}
C_{(a,c)} v_{T^{(a,c)}}$ with all four $C_{(a,c)} \ne 0$.  Consider
the coefficient of $v_{T^{(q, q+2)}}$ in $T_q b$.  This SYT is
\emph{not} in $\mathcal{S}_n$, so it is not in the support of $b$;
hence the coefficient of $v_{T^{(q, q+2)}}$ in $-b$ is $0$.  On
the other hand, by Lemma~\ref{lem:j=q-offdiag},
$T_q v_{T^*}$ produces the term $\beta v_{T^{(q,q+2)}}$ with
$\beta \ne 0$.  We need to check that no other summand of $T_q b$
cancels this contribution.

A summand $T_q (C_{(a,c)} v_{T^{(a,c)}})$ produces a term proportional
to $v_{T^{(q, q+2)}}$ only if either
\begin{itemize}[topsep=0.3em,itemsep=0.1em]
\item $T^{(a,c)} = T^{(q, q+2)}$ (impossible, since
  $T^{(q,q+2)} \notin \mathcal{S}_n$),
\item or $s_q T^{(a,c)} = T^{(q, q+2)}$.
\end{itemize}
$s_q T^{(a,c)}$ swaps the cells of $q$ and $q+1$ in $T^{(a,c)}$.
The resulting SYT has the same column-$2$ entries as $T^{(a,c)}$
except that any occurrence of $q$ or $q+1$ in column $2$ is swapped.

Among the four $(a,c) \in \mathcal{S}_n$:
\begin{itemize}[topsep=0.3em,itemsep=0.1em]
\item $(a,c) = (q+1, q+2)$: column~$2$ entries are $q+1, q+2$, both in
  column $2$.  Swap of $q,q+1$ moves $q+1$ from $(1,2)$ to $(q,1)$ and
  $q$ from $(q,1)$ to $(1,2)$ (provided $T^{(a,c)}(q,1) = q$, which
  Lemma~\ref{lem:top-col1} gives).  Resulting SYT: column-$2$ entries
  $q, q+2$, i.e., $T^{(q, q+2)}$.  \emph{This is the only contributor.}
\item $(a,c) = (q+1, q+3)$: same as above for the position of $q+1$
  at $(1,2)$.  Swap $q \leftrightarrow q+1$: now $q$ at $(1,2)$,
  $q+1$ at $(q,1)$.  Column-$2$ entries: $q, q+3$, i.e., $T^{(q, q+3)}$.
  \emph{Different from $T^{(q, q+2)}$.}
\item $(a,c) = (q+2, q+4)$: here $a = q+2 \ge n-2$, so by
  Lemma~\ref{lem:top-col1} the identity prefix has length $a - 1 = q+1$,
  giving $T^{(a,c)}(q, 1) = q$ and $T^{(a,c)}(q+1, 1) = q+1$.  Hence
  $q$ and $q+1$ are in column~$1$ at adjacent rows, so
  $T_q v_{T^{(a,c)}} = -v_{T^{(a,c)}}$ (same-column case).  \emph{No
  off-diagonal contribution.}
\item $(a,c) = (q+3, q+4)$: here $a = q + 3 = n - 1$, identity prefix
  length $q + 2$, so analogously $T_q v_{T^{(a,c)}} = -v_{T^{(a,c)}}$.
\end{itemize}

Therefore the coefficient of $v_{T^{(q, q+2)}}$ in $T_q b$ equals
$C_{(q+1, q+2)} \beta \ne 0$, while the same coefficient in $-b$ is
$0$.  Hence $T_q b \ne -b$, so $b \notin E_q^-$.
\qed
\end{proof}

\begin{proposition}[$B \not\subseteq E_j^-$ for $q+1 \le j \le n-1$]\label{prop:not-Ej-minus-higher}
For each $j \in \{q+1, q+2, \ldots, n-1\}$, $T_j b \ne -b$.
\end{proposition}

\begin{proof}
We treat the three sub-ranges separately.

\paragraph{$j = q + 1 = \ell - 1 = n - 3$.}  Positions of $j, j+1$ in
each support SYT $T^{(a,c)}$:
\begin{itemize}[topsep=0.3em,itemsep=0.1em]
\item $(a,c) = (n-3,n-2)$: $j = a$ at $(1,2)$, $j+1 = c$ at $(2,2)$.
  \emph{Same column}: $T_j v = -v$.
\item $(a,c) = (n-3,n-1)$: $j = a$ at $(1,2)$, $j+1 = n-2$ in column~$1$
  (specifically at row $n-3$).  \emph{Off-diagonal}, axial distance
  $1 - (-(n-4)) = n-3$.
\item $(a,c) = (n-2,n)$: $j = n-3$ in column~$1$ at row $n-3$,
  $j+1 = a = n-2$ at $(1,2)$.  \emph{Off-diagonal}, axial distance
  $-(n-3)$.
\item $(a,c) = (n-1,n)$: $j, j+1 = n-3, n-2$ both in column~$1$ at
  rows $n-3, n-2$.  \emph{Same column}: $T_j v = -v$.
\end{itemize}
Swap of $j, j+1$ in $T^{(n-3, n-1)}$: $j = n-3$ moves from $(1, 2)$
to column~$1$ row $n-3$; $j+1 = n-2$ moves from column~$1$ row $n-3$
to $(1, 2)$.  Resulting SYT: column~$2$ entries $(n-2, n-1)$, i.e.\
$T^{(n-2, n-1)}$, which is outside $\mathcal{S}_n$.
Swap of $j, j+1$ in $T^{(n-2, n)}$: analogously yields $T^{(n-3, n)}$,
also outside $\mathcal{S}_n$.  Since these two swap-targets are
distinct from each other and both have $|d| = n-3 \ge 3$
(non-degenerate Hoefsmit block), each contributes a nonzero
coefficient to $T_{q+1} b$ on a SYT outside $\mathcal{S}_n$, while
$-b$ is supported entirely in $\mathcal{S}_n$.  Hence $T_{q+1} b \ne -b$.

\paragraph{$j = \ell = n-2$.}  By Theorem~\ref{thm:T-ell}, $T_\ell b = q b$.
For generic $q \ne -1$, $qb \ne -b$, so $b \notin E_\ell^-$.

\paragraph{$j = \ell+1 = n - 1$.}  Positions of $j, j+1 = n-1, n$ in
each support SYT:
\begin{itemize}[topsep=0.3em,itemsep=0.1em]
\item $(n-3, n-2)$: column~$1$ entries are $\{1,2,\ldots,n-4,n-1,n\}$
  placed at rows $1,\ldots,n-2$ in increasing order, so $n-1$ at row
  $n-3$ and $n$ at row $n-2$ of column~$1$.  \emph{Same column}:
  $T_{n-1}v = -v$.
\item $(n-3, n-1)$: column~$1$ omits $n-3$ and $n-1$, so column~$1$ =
  $\{1,\ldots,n-4,n-2,n\}$ at rows $1,\ldots,n-2$.  $n-1$ at $(2,2)$,
  $n$ at row $n-2$ of column~$1$.  Different row, column.
  Off-diagonal.
\item $(n-2, n)$: column~$1$ = $\{1,\ldots,n-3,n-1\}$.  $n-1$ at row
  $n-2$ of column~$1$, $n = c$ at $(2,2)$.  Off-diagonal.
\item $(n-1, n)$: $n-1, n$ at $(1,2), (2,2)$.  Same column~$2$.
  $T_{n-1} v = -v$.
\end{itemize}
Swap of $j, j+1$ in $T^{(n-3, n-1)}$: $n-1$ moves from $(2, 2)$ to
column~$1$ row $n-2$; $n$ moves from column~$1$ row $n-2$ to $(2, 2)$.
Resulting SYT: column~$2$ entries $(n-3, n)$, i.e.\ $T^{(n-3, n)}$,
outside $\mathcal{S}_n$.
Swap in $T^{(n-2, n)}$: analogously yields $T^{(n-2, n-1)}$, also
outside $\mathcal{S}_n$.  Same argument as in the $j = q+1$ case:
each off-diagonal contribution gives a nonzero coefficient on a
distinct SYT outside the support of $b$.  Hence $T_{n-1} b \ne -b$.
\qed
\end{proof}

\subsection{$E_j^+$ non-containment for $j \ne \ell$}

\begin{proposition}[$B \not\subseteq E_j^+$ for $j \ne \ell$]\label{prop:not-Ej-plus}
For $j \in \{1, \ldots, n-1\} \setminus \{\ell\}$, $T_j b \ne q b$.
\end{proposition}

\begin{proof}
For $j \le q - 1$, $T_j b = -b$ by Proposition~\ref{prop:Ej-minus},
which equals $qb$ iff $-1 = q$, impossible for generic $q$.  For
$j \ge q$, the same off-diagonal arguments of
Propositions~\ref{prop:not-Eq-minus}--\ref{prop:not-Ej-minus-higher}
show that $T_j b$ has a nonzero component on an SYT outside
$\mathcal{S}_n$; but $qb$ is supported on $\mathcal{S}_n$.
Hence $T_j b \ne qb$.
\qed
\end{proof}

This completes the proof of Theorem~\ref{thm:main}.

\section{Computational verification}\label{sec:verify}

\begin{theorem}[Verification]\label{thm:verify}
Theorem~\ref{thm:main}, together with the support description of
Theorem~\ref{thm:may13-support}, was verified
computationally at $q = 2, 3, 4, 5, 6, 7$ (i.e., $n = 6, \ldots, 11$)
mod $P = 100\,003$ at $q = 11$.  Across these six cases:
\begin{itemize}[topsep=0.3em,itemsep=0.1em]
\item $\dim B = 1$ in all cases.
\item Support of $b$ matches $\mathcal{S}_n$ exactly, all four
  coefficients nonzero.
\item $T_j b = -b$ for $1 \le j \le q-1$ in all cases.
\item $T_q b \ne -b$ in all cases (verified by recovering a nonzero
  coefficient on $T^{(q, q+2)}$).
\item $T_\ell b = q b$ in all cases (Theorem~\ref{thm:T-ell}).
\end{itemize}
\end{theorem}

\begin{proof}
Scripts
\texttt{\textasciitilde/projects/scratch/2026-05-14-extended-sign-kill-proof/verify\_support\_argument.py}
and \texttt{check\_S1Omega\_zero.py}.  The scripts construct the
Hoefsmit $T_i$ matrices for $\lambda = (2,2,1^q)$, build $\Om$
explicitly, and check all of the above by direct linear algebra.
\qed
\end{proof}

\section{Discussion}\label{sec:discussion}

\subsection{What is new}

The extended sign-kill conjecture (May-14 evening) asserts, for
$\lambda = (\hatl, 1^q)$ with $\ell(\hatl) \ge 2$,
\[
   B \subseteq E_j^-\!\mid_{\Vlam} \;\Longleftrightarrow\; j \le \tau(\lambda),
\]
where $\tau(\lambda) = q + 3 - \hatl_1 - \hatl_2$.  For $\hatl = (2,2)$,
$\tau = q - 1$.

The May-13 paper \texttt{2026-05-13-non-hook-22-1n.tex} proves the
$j = 1$ case (containment in $E_1^-$).  The May-20 paper proves the
analogous $j = 1$ case for $\hatl = (3,2)$ at $q \ge 3$.  Neither
paper addresses $j \ge 2$.

The present paper closes the $\hatl = (2,2)$ row of the conjecture
in both directions: $B \subseteq E_j^-$ for $j \le q-1$, and
$B \not\subseteq E_j^-$ for $j \ge q$.  The proof is short and
extracted directly from the support structure already in May-13;
the contribution is the observation that the support immediately
forces $T(i,1) = i$ for $i \le n-4$, plus the sharpness argument
using the non-vanishing of the four support coefficients.

\subsection{The general support rule}\label{sec:support-rule}

The argument generalizes mechanically.  Define, for any
$H_q(S_n)$-submodule $W \subseteq V_\lambda$,
\[
   \mathrm{supp}(W) \;:=\; \{ T \in \mathrm{SYT}(\lambda) :
   \text{some } w \in W \text{ has nonzero coefficient on } v_T \},
\]
i.e., the union of the seminormal supports of all vectors in $W$.

\begin{proposition}[Support rule]\label{prop:support-rule}
Suppose $W \subseteq V_\lambda$.  If every $T \in \mathrm{supp}(W)$
has $j$ and $j+1$ in the same column of $T$, then
$W \subseteq E_j^-\mid_{V_\lambda}$.
\end{proposition}

\begin{proof}
Let $w \in W$, $w = \sum_{T \in \mathrm{supp}(W)} c_T v_T$ (with
some $c_T$ possibly zero).  By the Hoefsmit same-column rule,
$T_j v_T = -v_T$ for each $T$ in the support of $w$.  Linearity
gives $T_j w = -w$. \qed
\end{proof}

This is the crux of the proof for $\hatl = (2, 2)$ above (where the
support has $T(i,1) = i$ for $i \le n-4$, so $j, j+1$ are
column-$1$-adjacent for $j \le n - 5 = q - 1$).

\subsection{Empirical extension to $\hatl = (3, 2)$}

Computational verification (mod $P = 100\,003$ at $q = 11$) for
$\lambda = (3, 2, 1^q)$ at $q = 3, 4, 5, 6$:

\begin{itemize}[topsep=0.3em,itemsep=0.1em]
\item $\dim B = 2$ in all four cases, and $|\mathrm{supp}(B)| = 26$
  in all four cases.
\item Every $T \in \mathrm{supp}(B)$ has column-$1$ identity prefix
  of length $\ge q - 1$, i.e., $T(i, 1) = i$ for $i = 1, \ldots, q-1$.
  Witnesses: SYTs with $T(q, 1) \ne q$ (specifically $T(q, 1) = q+1$)
  achieve the minimum prefix $q - 1$.
\item Hence by Proposition~\ref{prop:support-rule},
  $B \subseteq E_j^-$ for $j \le q - 2$ — exactly the predicted
  threshold $\tau = q + 3 - 3 - 2 = q - 2$.
\item Sharpness ($B \not\subseteq E_j^-$ for $j \ge q-1$) is also
  computationally verified: the off-diagonal swap-targets at $j = q-1$
  produce SYTs outside the support, by the same argument as
  Section~\ref{sec:main}.
\end{itemize}

The structural proof would require characterizing $\mathrm{supp}(B)$
explicitly for $(3, 2, 1^q)$.  The May-20 paper constructs $B$ as the
image of the three-branch decomposition under
$R'_{\ell+1}\cdots R'_{n-1}$, and the resulting support is the
26-element set above.  Tracing the construction to extract the
support description would close the $\hatl = (3, 2)$ row of the
conjecture.  We leave this as future work.

\subsection{Extension blueprint for general $\hatl$}

The pattern that emerges:

\begin{conjecture}[Support structure conjecture]\label{conj:support}
For $\lambda = (\hatl, 1^q) \vdash n$ with $\ell(\hatl) \ge 2$, the
seminormal support $\mathrm{supp}(B^{(\lambda)}_{\ell+1} V_\lambda)$
consists exactly of those SYTs $T$ of $\lambda$ satisfying
$T(i, 1) = i$ for $i \le q + 3 - \hatl_1 - \hatl_2$.
\end{conjecture}

Conjecture~\ref{conj:support}, combined with
Proposition~\ref{prop:support-rule}, would prove the full extended
sign-kill conjecture (the ``$\subseteq$'' direction).  Sharpness in
each case would follow from witness SYTs at the boundary prefix
length, analogously to Section~\ref{sec:main}.

The top-two-rows independence $\tau = q + 3 - \hatl_1 - \hatl_2$
would then have a clean interpretation: it is the minimum
identity-prefix length in the support of $B$, and this minimum is
determined by the first two rows of the column-$1$-deletion shape
$\hatl^{(\ge 2)}$ (with $|\hatl^{(\ge 2)}_{\le 2}| = \hatl_1 + \hatl_2 - 2$).

\subsection{Containment versus identity}

Theorem~\ref{thm:main} gives $B \subseteq E_\ell^+ \cap \bigcap_{j=1}^{q-1}
E_j^-$.  The intersection on the right is computationally much larger
than $\dim B = 1$ (e.g., dimension $2$ for $\lambda = (2,2,1^q)$ at all
tested $q \ge 2$, see May-14 evening paper's Table~2).  So the
eigenspace constraints do not pin down $B$ exactly; the closed-form
$\dim B = f^{\lambda^{(\ge 2)}} = 1$ is a sharper fact, capturing structure not
detected by $T_j$-eigenspaces alone.

\subsection{Status of the broader iso programme}

This result closes one row of the May-14 evening conjecture table
but does not affect the broader iso programme (the May-13--May-14
attempts to identify $B$ as a known representation-theoretic
object).  Those attempts were closed off by the May-14 morning and
afternoon refutations.  The present contribution is, in effect, a
``cleaning up'' of the extended sign-kill empirical observation
into a structurally established theorem at the simplest non-hook
shape.

\section{Status}\label{sec:status}

\paragraph{Established.}
\begin{itemize}[topsep=0.3em,itemsep=0.1em]
\item Theorem~\ref{thm:main}: complete characterization of $T_j$-
  eigenspace containments of $B$ for $\lambda = (2,2,1^q)$, $q \ge 2$.
\item Sharpness for both $E_j^-$ and $E_j^+$ containments.
\end{itemize}

\paragraph{Open.}
\begin{itemize}[topsep=0.3em,itemsep=0.1em]
\item Extended sign-kill at $\hatl = (3, 2)$ at $j \ge 2$ (requires
  May-20-style support description; computationally consistent with
  threshold $\tau = q - 2$).
\item Extended sign-kill at general $\hatl$ (requires support
  description for general multi-corner shapes).
\item Structural mechanism for the top-two-rows independence
  $\tau = q + 3 - \hatl_1 - \hatl_2$ (independent of lower rows of
  $\hatl$).
\end{itemize}

\paragraph{Push status.}
TBD --- depends on PAT refresh.  Prior unpushed-commit count
on \texttt{clio-vega/proofs}: 64 (May-14 evening session).

\paragraph{Scripts.}
All computations:
\texttt{\textasciitilde/projects/scratch/2026-05-14-extended-sign-kill-proof/}.
\begin{itemize}[topsep=0.3em,itemsep=0.1em]
\item \texttt{explore\_22\_family.py} --- ranks and intersections of
  $E_j^\pm \cap B$ for $\lambda = (2,2,1^q)$ at $q = 2,3,4$.
\item \texttt{identify\_22.py} --- explicit basis vector of $B$ in
  Hoefsmit basis, support count, normalized coefficients.
\item \texttt{check\_S1Omega\_zero.py} --- verifies $S_1\Om = 0$ as a
  matrix on $\Vlam$ for $q \ge 2$ in the $(2,2,1^*)$ family, and
  diagnostic checks on other families ($(3,2,1^*)$, etc.).
\item \texttt{verify\_support\_argument.py} --- end-to-end verification
  of Theorem~\ref{thm:main}: support matches $\mathcal{S}_n$, the
  column-$1$ pattern, and the $T_j$-action at $j \le q-1$ and $j = q$.
\end{itemize}

\end{document}
